% Une hypertable : des tranches par jour ; une tranche compressée range, pour
% chaque capteur, des colonnes de valeurs ; la rétention supprime une tranche.
\documentclass[border=6pt]{standalone}
\usepackage{sd}
\begin{document}
\begin{tikzpicture}[x=1mm, y=1mm, note/.append style={font=\sffamily\small},
    t/.style={draw=#1, fill=#1Bg, line width=0.7pt, minimum width=17mm, minimum height=9mm, font=\sffamily\footnotesize, align=center, inner sep=1pt},
    col/.style={draw=Dslate, fill=DslateBg, line width=0.5pt, anchor=west, minimum height=5mm, font=\ttfamily\footnotesize, inner sep=1.5pt}]
  \node[leg, anchor=west] at (0,26) {hypertable, une tranche par jour};
  \foreach \i/\j/\c in {0/{1\textsuperscript{er} janv.}/Dred, 1/2 janv./Dblue, 2/3 janv./Dblue, 3/…/Dblue, 4/7 janv./Dgreen} {
    \node[t=\c] (j\i) at (10+19*\i,14) {\j};
  }
  \node[note, anchor=north, text=Dred, align=center] at (10,7) {drop\_chunks :\\fichier supprimé};
  \node[note, anchor=north, text=Dgreen, align=center] at (86,7) {tranche courante,\\en lignes};
  \node[note, anchor=north, align=center] at (48,7) {tranches compressées};
  % détail d'une tranche compressée
  \begin{scope}[xshift=112mm]
    \node[leg, anchor=west] at (0,24) {une tranche compressée, capteur 42};
    \node[col, minimum width=58mm] at (0,16) {ts : 00:00:00 +10 +10 +10 …};
    \node[col, minimum width=58mm] at (0,9) {temp. : 20,01 20,03 20,02 …};
    \node[col, minimum width=58mm] at (0,2) {hum. : 51,20 51,22 51,19 …};
    \node[note, anchor=west, align=left] at (60,16) {différences de\\différences : 0, 0, 0};
    \node[note, anchor=west, align=left] at (60,5.5) {valeurs voisines :\\XOR, peu de bits};
  \end{scope}
  \draw[arr=Dslate, dash pattern=on 2pt off 1.5pt] (j2.north) |- (110,21);
\end{tikzpicture}
\end{document}
